
Two neural networks can compute the same function yet occupy nearly orthogonal positions in weight space --- a geometric reality with concrete consequences for any method that learns from raw neural network weights. This paper provides the first empirical quantification of scaling and sign-flip symmetry orbit diameters as cosine distances in a real MLP model zoo, finding that both symmetry types produce geometrically large orbits (mean cosine distances 0.32 and 1.07, respectively; 100\% of oracle-constructed pairs exceed the 0.05 significance threshold across N=2,500 total orbit pairs). A probe of the Neural Functional Transformer (NFT), a permutation-equivariant encoder, reveals an asymmetric result: NFT is measurably non-invariant to scaling orbits (within-orbit vs. cross-orbit embedding gap +0.024, 95\% CI=[0.0232, 0.0245], entirely above zero), while NFT exhibits significant approximate invariance to sign-flip orbits by construction (gap=-0.0007, 95\% CI=[-0.0008, -0.0005], entirely below zero). An audit of the standard majority-sign canonicalization algorithm reveals that it fails to produce a unique canonical form for 85.6\% of Schürholt MNIST zoo models, due to an arithmetic property of even input dimension (d\_in=784) that the weight space learning literature has not previously characterized. Scaling canonicalization increases PCA explained variance ratio from 0.055 to 0.086 at k=20 principal components, confirming geometric redundancy is removed, but linear probe $R^2$ and NFT Spearman $\rho$ are not improved at the available zoo scale (N=500); all property prediction comparisons are statistically underpowered (95\% CI width $\approx 0.6)$. These findings establish a measurement-grounded framework for symmetry canonicalization in weight space learning: specifying which symmetries are geometrically large, which a given encoder already handles, which standard algorithms fail and why, and what experimental scale is required to detect downstream improvement.

